\section{Developer Experience, Sustainability, and Roadmap}
\label{sec:sustainability}

\subsection{Developer Experience and Infrastructure}
\label{sec:devex}

The development infrastructure is designed around one constraint: a project of
this scope spans many components with cross-cutting concerns, and the tooling
must make that manageable without forcing artificial coupling between parts
that should evolve independently.

\paragraph{Repository structure.}
The project uses a multi-repo layout under a single GitHub organization. Each
component has its own repository, its own release cycle, and its own CI
pipeline. Components can be updated independently as long as they do not
introduce breaking changes to the interfaces they expose. A breaking interface
change requires a coordinated update across all affected components, enforced
by the versioning mechanism described below.

\begin{lstlisting}[language=bash, caption={Repository structure}]
org/
  blueprint       <- architecture documents and this document
  kernel-layer    <- eBPF programs, Event Bus daemon
  knowledge       <- Kuzu integration, graph schema, Graph Daemon
  ai-layer        <- provider abstraction, MCP client
  compositor      <- Smithay-based Wayland compositor (Rust)
  desktop-shell   <- taskbar, launcher, shell (Tauri + TypeScript)
  ui-kit          <- shared component library (TypeScript + Tailwind)
  os-sdk          <- shared Rust crate for system integration
  app-files       <- File Manager
  app-settings    <- Settings
  app-terminal    <- Terminal
  app-store       <- Store
  module-sdk      <- SDK for shell module development
  themes          <- GTK4, Wine .msstyles, token source
  distro          <- build system, ISO generation, meta-packages
\end{lstlisting}

\paragraph{Build system.}
\emph{mkosi}, maintained by the systemd team, is the target tool for image
generation. It builds bootable OS images from a declarative configuration,
specifying which packages to include, which files to add, and how to configure
the system. It produces UEFI-bootable images directly and supports multiple
output formats: ISO, disk image, and container. It integrates cleanly with the
Open Build Service. Confirming that mkosi handles the full build pipeline in
practice is an open item pending actual use.

The \emph{Open Build Service} handles packaging for all custom components.
OBS builds RPM packages across architectures for both candidate base
distributions, signs packages with the project GPG key, and hosts the package
repository that end-user systems pull from. Both Fedora and OpenSUSE targets
can be built from the same spec files, which is one practical reason both
remain viable candidates.

\paragraph{Update system.}
Updates arrive from two independent streams managed separately. The base
system, meaning the kernel and base distribution packages, is handled by the
native package manager (dnf or zypper depending on the chosen base), with
Snapper automatically taking a Btrfs snapshot before every update. Project
components, meaning the compositor, shell, event bus, graph daemon, AI layer,
and first-party applications, are distributed via the project's own OBS
repository as RPM packages managed as a coordinated meta-package.

The update behavior is designed to be non-intrusive. Security updates for both
streams are downloaded automatically in the background and applied at the next
shutdown or reboot, never during an active session. No forced reboots, no
countdown timers. All other updates are downloaded silently in the background;
a single calm notification appears when they are ready. The user applies them
when convenient. In managed environments, administrators can configure mandatory
update windows; even then, no mid-session interruption occurs.

Rollback is available for every update: Btrfs and Snapper record a snapshot
before each change, and the user can boot into a previous snapshot from the
boot menu and roll back the entire system without any additional tooling.

Project components communicate over defined IPC interfaces. A breaking
interface change increments the major version number and must be shipped as a
coordinated meta-package update, so that incompatible component versions cannot
coexist. Components verify interface compatibility at startup:

\begin{lstlisting}[language=Rust, caption={Interface compatibility check at startup}]
fn check_compatibility() -> Result<()> {
    let daemon_version = get_graph_daemon_version()?;
    if daemon_version.major != EXPECTED_MAJOR {
        bail!("Graph Daemon interface version mismatch. Run system update.");
    }
    Ok(())
}
\end{lstlisting}

A component that detects a version mismatch refuses to start and surfaces a
clear message. The only resolution is applying the pending system update.

\paragraph{CI/CD pipeline.}
Tests run at three levels of integration. Table~\ref{tab:ci-triggers} shows
when each level runs.

\begin{table}[htbp]
\caption{CI trigger summary}
\label{tab:ci-triggers}
\begin{tabularx}{\linewidth}{lX}
\toprule
\textbf{Trigger} & \textbf{What runs} \\
\midrule
Every push or pull request     & Unit tests for the affected repository \\
Pull request merge to main     & Unit tests, lint, and build check \\
Daily at 02:00                 & Cross-component integration tests against all main branches \\
Release candidate              & Full VM tests and smoke tests \\
\texttt{kernel-layer} push     & Unit tests plus VM-based eBPF tests \\
\bottomrule
\end{tabularx}
\end{table}

Unit tests run per repository against mocks from \texttt{os-sdk}. Rust
components use \texttt{cargo test}; TypeScript components use Vitest. No PR
merges without green unit tests. The \texttt{kernel-layer} is a special case:
eBPF logic that can be tested without a real kernel runs on the host as normal
Rust unit tests, while eBPF programs themselves run in a minimal QEMU instance
that CI spins up automatically.

Cross-component integration tests verify that two or more components
communicate correctly over their real IPC interfaces. They run on the host
without a full system boot, starting only the directly involved daemons as real
processes against real IPC sockets. A failing integration test with all unit
tests passing means interface drift between repositories, which is the primary
failure mode this level exists to catch. Examples: the Event Bus emits a
\texttt{FileAccessEvent} and the Graph Daemon writes correctly to Kuzu; a theme
token change in Settings propagates to all running Tauri processes; the shell
sends a keybinding action and the Input Manager calls the correct handler.

Full VM tests boot a complete system image in QEMU and run smoke tests: system
boots without crash; Spotlight opens, finds an application, and launches it;
eBPF events arrive and land in the Knowledge Graph; theme switch propagates
correctly; module install and uninstall behave as expected. These tests are
slow and run on release candidates only. The VM image is built by the same
script used for real releases; no separate test image exists.

\paragraph{Local development workflow.}
Development happens at three levels of integration, each suited to a different
class of work. The common case is developing a single component in isolation.
Each repository has its own development setup with \texttt{cargo run},
\texttt{tauri dev}, or hot-reload as appropriate. Components that use the mocks
from \texttt{os-sdk} need no external dependencies running. Fast feedback loop,
runs entirely on the developer's machine.

For shell and compositor work, Smithay supports nested mode: the compositor
runs as a regular Wayland window inside the developer's existing desktop. The
full shell appears inside that window. No VM or extra hardware is needed. This
covers the majority of shell and compositor development; eBPF and kernel-level
components are not active in nested mode.

Full system work requires QEMU. This covers eBPF development, the boot
sequence, the update system, and anything that requires a real kernel or a
complete system environment. The \texttt{distro} repository contains a script
set for this workflow, including \texttt{build-image.sh} for a complete
bootable image, \texttt{run-vm.sh} for an interactive session,
\texttt{run-vm.sh --share \textasciitilde/projects/kernel-layer} to mount a
host directory via virtiofs into the VM so that the developer can compile on
the host and test in the VM without a full image rebuild, and snapshot
management scripts for saving and restoring VM state.

A strict rule applies to eBPF development: eBPF programs are developed and
tested exclusively inside the VM, never directly on the development machine.
A bad eBPF program can destabilize the kernel; the isolation of the VM is not
optional.

\paragraph{Mock strategy.}
Every IPC interface in the system has an official mock implementation that
lives in \texttt{os-sdk} directly alongside the interface definition. Mock and
real implementation share the same type definitions. An outdated mock is a
build error, not a runtime surprise: if an interface changes, the mock must be
updated in the same pull request.

\begin{lstlisting}[language=Rust, caption={Conditional mock usage in tests}]
#[cfg(test)]
use os_sdk::graph_daemon::mock::GraphDaemonMock as GraphDaemon;

#[cfg(not(test))]
use os_sdk::graph_daemon::client::GraphDaemonClient as GraphDaemon;
\end{lstlisting}

There are no per-repository custom stubs and no fixture files that can drift
from the real schema. The mock is the contract.

\paragraph{Documentation.}
Two documentation sites are maintained, both self-hosted. The user-facing wiki
at \texttt{wiki.[project].dev} is written by hand in the spirit of the Arch
Wiki: prose that explains intent and reasoning, not only what a setting does.
It covers concepts, how-tos, troubleshooting, and architecture explanations.
Community contributions are accepted via pull requests reviewed by the core
team. The wiki source lives in the \texttt{blueprint} repository alongside the
architecture documents.

The developer API reference at \texttt{api.[project].dev} is auto-generated
from code comments: \texttt{rustdoc} for Rust components,
\texttt{typedoc} for TypeScript components, aggregated and hosted under a single
domain so there is one place to search across all components. Keeping it
current is a CI gate: a pull request that removes or modifies a public API
without updating its documentation comment does not merge.

\subsection{Sustainability and Business Model}
\label{sec:business}

This is an open source project. The code is and remains open; that is a
foundational commitment, not a marketing position. At the same time, building
and maintaining a full operating system is a serious engineering effort that
requires sustained funding. The goal is not to extract value from users but to
find revenue streams that are structurally aligned with the project's values:
user sovereignty, transparency, and European digital autonomy.

\paragraph{Principles.}
Four constraints rule out entire categories of business models upfront. Core
OS functionality is not feature-gated behind a paywall; the desktop, the
permission system, the Knowledge Graph, and the AI layer are open and available
to everyone. There is no advertising, no telemetry sold to third parties, and
no data brokering. Venture capital with growth-at-all-costs expectations is
incompatible with a project built on user trust and is therefore excluded.
Enterprise features are an acceptable source of revenue; making the core worse
for non-paying users is not.

\paragraph{Store revenue.}
The Store's tiered commission model, described in Section~\ref{sec:store},
generates revenue proportional to ecosystem health. The project earns
meaningfully only if developers earn well first. At early stage this is
negligible. At meaningful adoption it becomes a significant funding stream. The
Store handles application sales, theme sales, module distribution, and profile
package distribution under a single infrastructure.

\paragraph{Enterprise services.}
The Signed Profile Package system described in Section~\ref{sec:profiles} is
already designed with organizations in mind. The \texttt{.lenv} format is open
and documented; anyone can author packages. Organizations deploying managed
environments at scale will often want professional help: authoring and
maintaining custom packages, deployment tooling, update management, and support
SLAs. This is a natural B2B offering that does not touch the open source core.
The analogy is Red Hat selling support for software that anyone can download
for free.

\paragraph{OEM and hardware partnerships.}
At sufficient maturity and hardware compatibility, small European laptop
manufacturers may want to ship the system pre-installed. The structural
alignment is good: the manufacturer gets a differentiated product, users get a
well-tested hardware and software combination, and the project gets sustainable
revenue from certification and validation work and an ongoing support
relationship. This requires a stable release and a polished out-of-box
experience first and is a longer-term play. Relevant candidates in the European
market include Tuxedo Computers, the PINE64 ecosystem, and potential Framework
partnerships.

\paragraph{Support tiers.}
A support subscription aimed at power users, small businesses, and developers
who want guaranteed response times and direct access offers prioritized bug
reports, LTS release guarantees with defined support windows, and a direct
communication channel with the core team. No features are locked behind this;
it is a service layer on top of the open product, priced modestly enough to be
accessible to individual developers.

\paragraph{Public funding and grants.}
The European open source and digital sovereignty funding landscape is
underutilized by most projects. The Sovereign Tech Fund in Germany funds open
source infrastructure projects and explicitly values work that reduces
dependency on US technology monopolies. The NLnet Foundation funds internet and
open source projects with a privacy and user rights focus and has supported
numerous small to medium projects. The NGI (Next Generation Internet) programme
runs multiple EU funding calls per year for open source, privacy, and
decentralization work.

The European digital autonomy angle is not only a philosophical position; it is
a genuine funding argument. A European open source operating system built on
sovereignty principles is precisely what these programmes exist to support.
Grants do not provide recurring revenue but they can fund specific development
phases without creating investor obligations or governance compromises.

\paragraph{Path to incorporation.}
A company makes sense when there is external money to manage, not before.
Premature incorporation adds administrative overhead with no corresponding
benefit. A realistic sequence: the project gains traction and a community
forms; grant applications are filed; Store revenue starts, providing the first
external money flow; the first enterprise inquiry for deployment support
arrives as the first B2B signal. At that point incorporation becomes
appropriate, probably as a GmbH in Austria or Germany, or as a European
cooperative structure if community governance is a priority.

The company's role is to employ core contributors, manage the Store
infrastructure, and sell enterprise and support services. The open source
project remains governed independently of the company: contribution policy,
release decisions, and core architecture are not dictated by commercial
interests. This separation must be written down explicitly before the company
is incorporated, not negotiated afterward under time pressure.

\subsection{Roadmap}
\label{sec:roadmap}

The roadmap is intentionally sparse at this stage. Defining detailed phase
scope before the foundational components are proven would produce a plan that
reflects assumptions rather than reality. Phases 1 through the public release
will be scoped as Phase 0 concludes and the first components are built and
measured.

\paragraph{Phase 0 - Foundation (current).}
Architecture and design decisions are documented in this blueprint. The proof
of concept consists of the Event Bus and Knowledge Graph running as a daemon
and collecting basic system events. The base distribution decision between
Fedora and OpenSUSE Slowroll is pending and will be made before Phase 1 begins.
The Kuzu benchmark described in Section~\ref{sec:knowledge-graph} at 5,000
events per second is a Phase 0 exit criterion.

\paragraph{Phase 1 and beyond.}
Phase scope will be defined once Phase 0 is complete. The ordering of
subsequent phases follows the architectural dependency graph: kernel layer and
event infrastructure before the graph, the graph before the AI layer, the
compositor before the shell, the shell before Store and module infrastructure.
Public release comes after a stable, tested combination of all layers exists.
